% Claim proof file for row claim_3_23 -- "SigmaOpenPathsWalks".
% Auto-generated stub by scaffold/solve_chapter.py at row start. A manager
% agent dedicated to writing the TeX proof fills in the proof body below.
%
% This file is a sibling of `main.tex` inside `<subsection>/tex/`. The
% `subfiles` package lets it render both standalone and as part of
% `main.tex` via `\subfile{...}`.
%
% The statement is restated above the proof so the file is self-contained
% when rendered on its own. The proof must:
%   - Stay close to the lecture notes' proof if one exists (always search
%     the LN first for a `\begin{proof}` block immediately following the
%     claim; copy / lightly edit when present).
%   - Otherwise use the LN paradigm and cite prior claims by their `ref`.
%   - Be self-contained enough that a mathematician can follow it without
%     re-reading the LN.
%   - Have no `sorry`, `omitted`, or "obvious" shortcuts.

\documentclass[main]{subfiles}

\begin{document}
% ref: claim_3_23 (proof, with statement restated for readability)
% title: SigmaOpenPathsWalks
\def\rowref{claim\_3\_23}
\phantomsection\label{claim_3_23}

% --- statement (restated) ---------------------------------------------------
\begin{Note}[SigmaOpenPathsWalks]{Prp}{restateprpsigmaopens}\label{prp:sigma_opens}
  Let $G=(J,V,E,L)$ be a CDMG. For $C \subseteq J \cup V$, and $w_1, w_2 \in J \cup V$, the following are
  equivalent:
  \begin{enumerate}
      \item there exists a $C$-$\sigma$-open \emph{path} between $w_1$ and $w_2$ in $G$;
      \item there exists a $C$-$\sigma$-open \emph{walk} between $w_1$ and $w_2$ in $G$;
      \item there exists a $C$-$\sigma$-open \emph{walk} between $w_1$ and $w_2$ in $G$ such that all its colliders lie in $C$ (and not just in $\Anc^G(C)$).
  \end{enumerate}
\end{Note}

% --- proof ------------------------------------------------------------------
\begin{proof}
% TODO: write the proof body.
\end{proof}

\end{document}
